% Section: Governance Model
\section{Governance Contract Model}
\label{sec:governance-model}

% TODO: Expand this section.
%
% Key topics:
% - Formal definition of governance contracts
% - Contract composition and inheritance
% - Enforcement mechanisms
% - Contract violation handling
% - Integration with HITL checkpoints and milestone synchronization

A \textit{governance contract} in the Reflector framework is a formal specification
of the behavioral constraints and obligations imposed on an autonomous agent or
agent group operating within a bounded scope. Governance contracts make implicit
expectations explicit and machine-enforceable.

\subsection{Contract Components}

A governance contract $C$ consists of the following components:

\begin{itemize}
  \item \textbf{Scope}: The bounded execution context within which the contract applies
    (e.g., a specific module, feature branch, or milestone phase).
  \item \textbf{Obligations}: Actions the agent is required to perform (e.g., emit
    audit events, halt at milestone boundaries, request approval for out-of-scope changes).
  \item \textbf{Prohibitions}: Actions the agent is forbidden from taking without explicit
    human approval (e.g., modifying infrastructure configuration, altering governance
    invariants).
  \item \textbf{Permissions}: Actions the agent may take autonomously within the contract
    scope (e.g., generating code, running tests, creating pull requests).
  \item \textbf{Escalation policy}: The protocol triggered when the agent encounters a
    situation not covered by the contract.
\end{itemize}

\subsection{Contract Enforcement}

Contracts are enforced by the Reflector runtime, which intercepts agent actions and
validates them against the applicable contract before execution. Prohibited actions
trigger an automatic halt and governance checkpoint escalation.
